% !TeX root = ../../../main.tex
\section{Double-Profile Intersection Elastography}
\label{sec:theory}

To understand DoPIo imaging, consider that an acoustic radiation force (ARF) excitation induces displacement of sub-resolution scatterers. At the moment of excitation, scatterers located at the center of the ARF point spread function (PSF), where the ARF amplitude is highest, displace more than those near the edges of the PSF, where the amplitude is lower~\cite{sarvazyan_Shear_1998, ostrovsky_NonDissipative_2006,sarvazyan_Acoustic_2021}. Scatterers located outside the ARF PSF do not initially displace. This spatial gradient in displacement creates scatterer shearing, which propagates outward from the center of the ARF PSF over time~\cite{czernuszewicz_Experimental_2013}.

If these scatterer displacements are tracked using a beam with a PSF that is confocal with, and equally wide as the ARF push beam, the resulting A-line will reflect the weighted average displacement of all of the excited scatterers~\cite{walker_fundamental_1995,mcaleavey_Estimates_2003}. On the other hand, if a confocal tracking beam has a PSF that is narrower than the ARF excitation, then the resulting A-line will reflect only the average motion of the scatterers near the center of the ARF PSF. These scatterers respond immediately to the ARF excitation and then undergo elastic recovery. In contrast, if displacements are tracked using a confocal beam with a PSF that is wider than the ARF excitation, the resulting A-line will average together both the initially displaced scatterers and those that begin to move later due to the propagation of scatterer shearing. As a result, the displacement profiles that will be measured from the A-lines generated by these different tracking approaches will differ. 

Since all three displacement profiles encode information about the mechanical response of the same region of tissue to the same applied ARF, they can be compared to reveal the speed at which scatterer shearing propagates laterally. For example, because scatterer shearing propagates faster in stiffer media, the time at which the narrow and wide tracking beams report the same average scatterer displacement occurs earlier. To quantify this phenomenon, we define the time at which the two displacement profiles intersect as the “time-intersect”, denoted as $t_{\textnormal{int}}$. This time-intersect is related to shear elastic modulus through an empirically derived model. Figure~\ref{fig:mechanism} illustrates DoPIo's conceptual framework.


\begin{figure}[tbp]
    \includesvg[width=\linewidth]{parts/dopio/fig/mechanism.svg}
    \caption{Conceptual illustration of the beam sequence and signal processing employed in DoPIo elastography. From left to right: (1) An acoustic radiation force excitation is applied with a predetermined focal configuration. (2-4) Scatterer displacements are tracked over time using two confocal tracking beams with different lateral widths. (5) The time at which the resulting displacement profiles intersect, denoted as $t_{\textnormal{int}}$, reflects the rate of lateral scatterer shearing and is related to shear elastic modulus through an empirically derived model.}
    \label{fig:mechanism}
\end{figure}

We previously introduced DoPIo and demonstrated it for quantifying shear elastic modulus \emph{in silico}~\cite{yokoyama_Doubleprofile_2019,yokoyama_Bessel_2023,muhtadi_Double_2025} and experimentally in select imaging and processing pipelines~\cite{yokoyama_Blind_2020,yokoyama_Assessing_2021,yokoyama_Quantitative_2022}. This work builds on our previous work by identifying the ARF excitation and tracking beam focal configurations and sequences, as well as the empirical model, that yield the most accurate and precise shear elastic modulus estimates. The performance of these optimized DoPIo methods is then evaluated \emph{in silico}, \emph{in vitro}, and \emph{ex vivo}.
